\subsection{A proposed one-prime cuspidal repair of configuration exactness}

\begin{quote}
\textbf{Research draft, not yet approved for integration.}
The argument below concerns the maximal cuspidally pro-\(r\) quotient,
not the full prime-to-\(p\) fibre. It must be independently audited before
any unconditional separation conclusion is drawn.
\end{quote}

Let \(C\) be a smooth proper connected complex curve of genus at least two,
let \(P\in C\), and fix distinct primes \(p,r\). Use full profinite groups
\[
 F=\pi_1(C-\{P\}),\qquad B=\pi_1(C),\qquad
 T=\pi_1(\operatorname{Conf}_2(C)).
\]
Projection to the second coordinate gives the full profinite exact sequence
\[
                   1\longrightarrow F\longrightarrow T
                     \longrightarrow B\longrightarrow1.       \tag{CR.1}
\]
Only full profinite exactness is used here. Put
\[
 \Delta=B^{(p')},\quad F'=F^{(p')},\quad
 K'=\ker(F'\longrightarrow\Delta),\quad
 F_r=F'/\ker(K'\longrightarrow(K')^{(r)}),\quad K=(K')^{(r)}.
                                                               \tag{CR.2}
Thus \(1\to K\to F_r\to\Delta\to1\) is exact and \(K\) is pro-\(r\).
Let \(T_r\) be the maximal cuspidally pro-\(r\) quotient of \(T^{(p')}\)
relative to its natural surjection onto \(\Delta\times\Delta\).

\paragraph{Proposed conclusion.}
The natural map induces an exact sequence
\[
                   1\longrightarrow F_r\longrightarrow T_r
                     \longrightarrow\Delta\longrightarrow1.   \tag{CR.3}
\]
The point requiring scrutiny is injectivity on the left. The following proof
is intended to establish it without completing the base in a fibre-only
homotopy exactness statement.

\paragraph{1. The cusp Frattini module and the centre.}
Fix a consistently oriented puncture generator \(b\in F\). The kernel
\(K'\) is the closed normal subgroup generated by its image in \(F'\).
The mod-\(r\) Frattini quotient
\[
                  M=K/\Phi_r(K)
\]
is canonically, after choosing the puncture lift, the completed permutation
module
\[
                         M\simeq\mathbb F_r[[\Delta]].           \tag{CR.4}
\]
Here \(\Phi_r(K)=\overline{K^r[K,K]}\). One way to check (CR.4) is to
take a cofinal system of characteristic open normal subgroups
\(H\subset\Delta\), with corresponding proper covers \(C_H\to C\).
Let \(F_H\subset F'\) be their inverse images and \(S_H\) the inverse
image of \(P\). The puncture localization sequence in mod-\(r\) homology
is
\[
 0\longrightarrow
   \mathbb F_r[\Delta/H]/\mathbb F_r\!\sum_{x\in\Delta/H}x
 \longrightarrow H_1(F_H,\mathbb F_r)
 \longrightarrow H_1(C_H,\mathbb F_r)\longrightarrow0.           \tag{CR.5}
\]
Transitions are induced by the covering maps. The inverse limit of the
right-hand groups is zero: dually, each finite collection of mod-\(r\)
characters of \(H\) dies on a sufficiently small open subgroup, which
may be replaced by its characteristic core in \(\Delta\). The inverse
limit of the displayed sum lines is also zero, since a deeper cover can
be chosen with degree divisible by \(r\), and pushforward then sends its
sum vector to zero. The inverse systems are finite, hence satisfy the
Mittag--Leffler condition. Continuity of cohomology for the closed subgroup
\(K'=\bigcap_HF_H\), followed by finite duality, identifies
\(H_1(K',\mathbb F_r)\) with the inverse limit of the middle groups.
This proves (CR.4).

We also need that \(F_r\) is centre-free. The group \(F'\) is free
pro-prime-to-\(p\), and \(K'\) has \(r\)-cohomological dimension at
most one. Equivalently, this follows by continuity from the vanishing of
\(H^2(F_H,\mathbb F_r)\) for the open free pro-prime-to-\(p\) groups
\(F_H\). Inflation
\[
 H^2(K,\mathbb F_r)\longrightarrow H^2(K',\mathbb F_r)
\]
is injective: a central extension of the pro-\(r\) group \(K\) by
\(\mathbb F_r\), whose pullback to \(K'\) splits, has a splitting map
from \(K'\) into a pro-\(r\) group, and this map factors through
\(K=(K')^{(r)}\). Thus \(H^2(K,\mathbb F_r)=0\), so \(K\) is a
free pro-\(r\) group. Its rank is greater than one by (CR.4), and hence
its centre is trivial. Since the proper hyperbolic surface group
\(\Delta\) is centre-free, an element of the centre of \(F_r\) belongs
to \(K\) and centralizes \(K\); it is therefore trivial.

\paragraph{2. A finite-group automorphism lemma.}
Let \(Q\) be a finite group with normal \(r\)-subgroup \(N\). The
subgroup of \(\operatorname{Aut}(Q)\) inducing the identity on both
\(Q/N\) and \(N/\Phi_r(N)\) is an \(r\)-group. Its restriction to
\(N\) is an \(r\)-group by the Frattini automorphism theorem for finite
\(r\)-groups. The kernel of that restriction fixes \(N\) pointwise and
\(Q/N\) pointwise. For such an automorphism \(a\), the function
\[
                     q\longmapsto a(q)q^{-1}
\]
takes values in \(Z(N)\), is constant on cosets of \(N\), and is a
crossed homomorphism from \(Q/N\) into the finite abelian \(r\)-group
\(Z(N)\). This identifies the kernel with a subgroup of an \(r\)-group.

It follows that the closed subgroup
\[
 \mathcal I=\{a\in\operatorname{Aut}(F_r):
           a|_\Delta=1,\ a|_M=1\}                              \tag{CR.6}
\]
is pro-\(r\). Indeed \(F_r\) is topologically finitely generated, so
its automorphism group is profinite and detected by characteristic finite
quotients \(Q\). The image \(N\) of \(K\) in any such \(Q\) is an
\(r\)-group; the surjection \(M\to N/\Phi_r(N)\) shows that the
preceding finite lemma applies to the image of \(\mathcal I\).

\paragraph{3. The full configuration action is prime-to-\(p\).}
Conjugation in (CR.1) preserves the kernel defining \(F_r\), and hence
gives a continuous homomorphism
\[
                    \rho:T\longrightarrow\operatorname{Aut}(F_r).
                                                               \tag{CR.7}
Let \(D\to T\) be the image of the full geometric peripheral group
around the diagonal. Its projection to the second factor \(B\) is
surjective, and it centralizes the oriented inertia generator \(b\).
No injectivity of \(D\to T\) is needed here. Since \(D\to B\) is
surjective, every \(t\in T\) can be written \(t=f d\), with
\(f\in F\), \(d\in D\).

Write \(g\in\Delta\) for the common image of \(d\) under the two
proper-coordinate maps, and \(h_0\in\Delta\) for the image of \(f\).
The action of \(d\) on \(\Delta\) is conjugation by \(g\), and its
action on the cusp classes indexed by \(h\in\Delta\) is
\[
 [hbh^{-1}]\longmapsto[(ghg^{-1})b(gh^{-1}g^{-1})].
\]
The ambiguity in a chosen lift of \(h\) belongs to \(K\) and acts
trivially on \(M\), so this formula is well-defined. Inner conjugation
by \(f\) acts on the indices by left multiplication by \(h_0\).
Thus the combined action on \(M\) is
\[
              h\longmapsto g_1 h g_2^{-1},
       \qquad (g_1,g_2)=(h_0g,g),                              \tag{CR.8}
\]
and the action on \(\Delta\) is conjugation by \(g_1\). Both actions
therefore factor through the proper-coordinate map
\(\theta:T\twoheadrightarrow\Delta\times\Delta\).
For this factorization it is enough that the cusp classes generate \(M\);
the full isomorphism (CR.4) is not needed.

Consequently \(\rho(\ker\theta)\subset\mathcal I\). The image of
\(\rho\) is an extension of a pro-\(r\) group by a quotient of
\(\Delta\times\Delta\). It is therefore pro-prime-to-\(p\).
To retain the proper-coordinate map without any faithfulness assertion
about (CR.8), use the joint image
\[
        E=\operatorname{im}(\rho,\theta)
          \subset\operatorname{Aut}(F_r)\times(\Delta\times\Delta).
                                                               \tag{CR.9}
\]
This is a pro-prime-to-\(p\) quotient of \(T\), and its kernel over
\(\Delta\times\Delta\) is pro-\(r\), by (CR.6). It thus factors
through the defining universal quotient \(T_r\).

The restriction of (CR.9) to \(F\) factors through \(F_r\), and its
first coordinate is the inner action of \(F_r\) on itself. This action
is faithful by centre-freeness. On the other hand the map \(F'\to T_r\)
kills \(\ker(K'\to(K')^{(r)})\), since the relative kernel of
\(T_r\to\Delta\times\Delta\) is pro-\(r\). We therefore have
\[
                     F_r\longrightarrow T_r\longrightarrow E
\]
with injective composite. This proves the left injectivity in (CR.3).
Right exactness follows from (CR.1), since passage to maximal
prime-to-\(p\) quotients is right exact and the further cuspidal quotient
is taken inside the kernel of the proper-coordinate map.

\paragraph{4. Proposed arithmetic consequence and its remaining audit.}
Prime-to-\(p\) specialization for the smooth proper pair
\((C\times C,\operatorname{diag}C)\), compatibly with a punctured
fibre, would transfer (CR.3) to a curve in characteristic \(p\).
The quotient is characteristic for the arithmetic action, so it gives an
arithmetic extension over \(\Pi_C\). The inverse image of a rational
point section is then the corresponding actual one-point cuspidally
pro-\(r\) group.

One must also check cusp compatibility in this last assertion. The full
peripheral group around the diagonal is an extension of the proper group
by its central geometric inertia. After passage to the prime-to-\(p\)
and then the cuspidally pro-\(r\) quotient, its kernel over the proper
group is exactly the image of that inertia, by right exactness. This
image is \(\mathbb Z_r(1)\), because its image in the retained fibre
\(F_r\) is injective. Thus the peripheral subgroups of two equal
proper rational sections would be the same inverse image in this one
diagonal peripheral group.

If all these steps survive independent audit, equal proper point sections
would give the cusp-compatible isomorphism required by the separately
proved conditional separation criterion. Until that audit is complete,
this draft does not claim that A1 has been settled.
